\section{Discussion}
\label{sec:discussion}

\subsection{Finding}
\label{sec:discussion-finding}

Increasing disentanglement strength substantially changed representation geometry without improving distance-based reliability estimation, while the domain information such estimation would require remained recoverable from the same representations by a non-distance-based reader.

\subsection{Interpretation}
\label{sec:discussion-interpretation}

This study audited a simple assumption: changing representation geometry should change the behavior of distance-based reliability estimators. The first part of this assumption held---condition number moved by two orders of magnitude across the disentanglement ladder (Section~\ref{sec:results-geometry}). The second did not: no geometry metric tracked Mahalanobis AUROC, and AUROC itself showed no trend across the ladder (Section~\ref{sec:results-geometry-auroc}).

A null result admits two straightforward explanations: that Mahalanobis distance specifically is a poor instrument, or that disentanglement training erased the domain-discriminative information a reliability estimator would need. Neither holds. Cosine-to-centroid similarity discards the covariance structure Mahalanobis distance depends on; pooled $k$-nearest-neighbor distance discards the notion of a class centroid altogether; Energy \citep{liu2020energy} discards representation geometry altogether in favor of the classifier's own logits; ViM \citep{wang2022vim} combines a residual-subspace geometric term with those same logits; a per-class kernel density estimator replaces Mahalanobis's single-Gaussian assumption with a non-parametric one. Despite these substantially different scoring principles, seven of the eight scorers tested converged on the same 0.39--0.42 AUROC (Section~\ref{sec:results-scorer-comparison})---a failure shared across formulas this different reflects what they share, not a defect specific to any one. Nor was the relevant information absent: three probes with no access to the training objective recovered domain membership at 0.72--0.81 AUROC at every disentanglement level tested (Section~\ref{sec:results-probe}), well above what these scoring rules reported on the identical embeddings. The information was not gone. It was present, recoverable, and specifically unreachable by seven of the eight scoring rules tested.

The eighth scorer, Energy, is an exception we report plainly rather than explain away. Among the eight evaluated scorers, Energy exhibited a modest monotonic trend with increasing disentanglement (Jonckheere--Terpstra exact $p=0.032$, one-sided; its two-sided counterpart coincides with the Kendall $p=0.065$ by the $\tau_b$--$J$ identity of Section~\ref{sec:stats}, so the two tests are one piece of evidence read two ways, not two), whereas the remaining seven scorers showed no consistent trend (Section~\ref{sec:results-scorer-comparison}). This isolated observation should be interpreted cautiously: it is one result out of eight scorers tested without correction for multiple comparisons, it was not part of the analysis plan fixed before the experiments ran, its Kendall's $\tau$ (0.44, exact $p=0.065$) does not clear the $|\tau|\geq0.3$-and-significant threshold that plan specifies (Section~\ref{sec:stats}), and the effect size is modest. We treat it as a hypothesis-generating observation rather than as evidence against the main conclusion, which we accordingly scope to seven of the eight scorers tested rather than all eight. That Energy alone showed this trend suggests that accessibility may depend not only on representation geometry but also on the inductive biases of individual scoring functions---whether this reflects a genuine property of energy-based scoring or a dataset-specific effect is an open question, as is whether it persists under different pretraining paradigms or model architectures.

Taken together, these findings indicate that \textbf{information can remain decodable while becoming largely inaccessible to non-probing reliability estimators.} This is not a claim that any one of these methods is poorly designed, and not a claim that disentanglement training destroys information. It is a claim about a gap between two properties of a representation---whether information is present, and whether it is organized in a form a specific family of estimators can use---that the ladder design in this study makes it possible to separate for the first time in this setting. This claim is scoped to the training regime studied here---domain-adversarial disentanglement evaluated against its own adversarial target domain (Section~\ref{sec:discussion-limitations})---not to representation learning or domain shift in general.

Section~\ref{sec:results-distance-gap}'s supporting evidence (a modest, direction-inconsistent norm difference; a majority-class attraction of 8.7--10.5 percentage points above baseline) indicates that the raw distance gap between domains is real but is not fully accounted for by either mechanism individually; neither is proposed here as a complete account of \emph{why} the gap exists. Understanding the origin of this systematic below-chance behavior remains future work.

\subsection{Relation to literature}
\label{sec:discussion-literature}

Mahalanobis-distance OOD detection \citep{lee2018simple} and its non-parametric successors, including pooled $k$-nearest-neighbor scoring \citep{sun2022out}, are typically evaluated by benchmarking AUROC across datasets and architectures, not by auditing whether the geometric assumption connecting representation and reliability estimate survives a specific training intervention. Energy-based scoring \citep{liu2020energy} and Virtual-Logit Matching \citep{wang2022vim} were developed specifically to move beyond a purely geometric notion of reliability by incorporating the classifier's own logits, making them a natural extension of the audit rather than a different question. This study does not propose a competitor to any of these methods; all five distance-based and three non-distance-based scorers are used here as instruments for the audit, alongside a cosine-similarity variant and a kernel density estimator that isolate specific structural assumptions from one another. The finding that seven of the eight degrade identically under this training regime is a statement about a premise most of these methods share, not a comparative ranking of them.

Domain-adversarial and orthogonality-based disentanglement training \citep{ganin2015unsupervised} is one of several strategies proposed to reduce reliance on shortcut features in deep classifiers \citep{geirhos2020shortcut}. That literature evaluates disentanglement primarily by its effect on classification robustness and shortcut reliance; this study evaluates a different, downstream consequence---what disentanglement does to a reliability estimator built on top of the resulting representation---and finds that the two do not move together. This does not bear on whether disentanglement succeeds at its stated goal, only on whether a distance-based reliability estimator layered on top of it continues to function as it did before the intervention.

The distinction this study draws between the presence of information in a representation and a specific estimator's ability to use it is closely related to, and constrained by, the probing-classifier literature's own caution that a supervised probe's success does not establish that a model itself uses the decoded information in the way the probe accesses it \citep{hewitt2019designing}. This study makes a narrower claim than that literature warns against overreaching toward: the probes here are not offered as evidence about what the disentanglement objective computes internally, only as an independent measurement of what remains extractable from its output representation, without characterizing the mechanism by which that representation was produced.

More broadly, these findings suggest that evaluating learned representations only through downstream predictive performance may overlook an equally important question: whether the representation is organized in a form that downstream reliability estimators can actually exploit. Geometry, decodability, and distance accessibility therefore appear to be related but distinct properties, and improvement in one should not automatically be interpreted as improvement in the others.

For practitioners combining domain-adversarial disentanglement training with a distance-based reliability gate in a deployed pipeline, this audit implies a specific caution: such a gate may not reliably flag inputs from the domain the disentanglement objective was trained against, even where it performs as expected against unrelated distributions. Where feasible, reliability estimation in such pipelines should be validated directly against the domain the training procedure targets, rather than assumed to transfer from general model qualification.

\subsection{Limitations}
\label{sec:discussion-limitations}

PAD-UFES is not an arbitrary target domain in this design: it is the domain the disentanglement architecture's adversarial branch is trained against. Since PAD-UFES also serves as the adversarial domain during training, our conclusions should be interpreted as characterizing representations learned under this training regime rather than all domain shifts.

The comparison against a conventional, non-disentangled baseline (Section~\ref{sec:results-baseline}) differs from the disentanglement ladder in architecture, representation dimensionality, and training recipe simultaneously, and is drawn from a single checkpoint. It is reported descriptively and does not support any causal attribution of the ladder's findings to disentanglement specifically, as opposed to architecture or dimensionality; a matched, non-adversarial checkpoint at the same architecture and 16-dimensional representation would be needed to make that attribution, and no such checkpoint exists in the inventory available for this study. Attributing the central finding specifically to disentanglement, rather than to architecture or dimensionality, remains an important direction.

The ladder itself spans three ordinal treatment levels with unequal per-level seed counts ($n=13$ total, $n=9$ on the common-seed subset), and we state the resulting limits on detection quantitatively rather than as a qualitative caveat (Section~\ref{sec:power}; Supplementary Section~S8). At $n=13$ the smallest $|\tau|$ that can reach $p\leq0.05$ under the exact permutation null is $0.436$ for the continuous--continuous design used in Section~\ref{sec:results-geometry-auroc} and $0.473$ for the 5/5/3 ordered ladder used in Sections~\ref{sec:results-scorer-comparison}--\ref{sec:results-probe} ($0.609$ on the $n=9$ common-seed subset). Two consequences follow, and we state both explicitly. The pre-specified success criterion this study applies---$|\tau|\geq0.3$ \emph{and} significant---lies below all three of those thresholds and was therefore unsatisfiable before any data were collected: an association of exactly the pre-specified size could not have been declared significant at this design. And the power to detect an expected $\tau=0.3$ is $0.27$, $0.23$, and $0.14$ for the three designs respectively; 80\% power is first reached at $\tau=0.54$, $0.57$, and $0.69$. The null results in this study are consequently informative about large associations and close to uninformative about moderate ones.

The bootstrap intervals in Supplementary Section~S9 say the same thing from the data side. Of the twenty null associations for which a bootstrap interval is valid, nineteen have a 95\% interval wide enough to contain $|\tau|=0.3$; the sole exception is ViM ($\tau=0.076$, 95\% CI $[-0.268,0.270]$), which is the only null result here that excludes an association of the pre-specified magnitude. For the remaining nineteen the honest statement is that this study does not distinguish ``no association'' from ``a moderate association this design cannot resolve.'' We report this because it bounds the central claim: the dissociation documented in Sections~\ref{sec:results-geometry-auroc}--\ref{sec:results-probe} rests on how consistently the point estimates fall near zero across the twenty associations tested outside the condition-number effect---four further geometry metrics against $\lambda_{orth}$, five against Mahalanobis AUROC, eight scorers, and three probes---and not on any individual interval being tight. Eighteen of those twenty have $|\tau|<0.3$; the two that do not (Energy, $\tau=0.443$; cosine-to-centroid, $\tau=0.321$) are reported as such in Section~\ref{sec:results-scorer-comparison} and are not counted as null. These twenty tests are not mutually independent---the geometry metrics derive from a common fitted precision matrix and the three probes share one feature set---so their agreement is weaker evidence than twenty independent replications would be. A single one of these tests, read alone, would not support the claim.

Exact-permutation testing was used throughout in place of asymptotic approximations not verified to hold at this scale, and $p$-values are reported without family-wise correction across the multiple metrics and scorers tested, on the stated basis that conclusions rest on consistency of effect across independent analyses rather than on any single threshold crossing; a reader applying a stricter correction should note that the one large, load-bearing effect in this study (condition number's association with $\lambda_{orth}$, $\tau=0.840$---which is the maximum value the 5/5/3 design permits, i.e.\ a perfectly monotone ladder) survives any standard correction, while every null result was already reported as non-significant before any correction was considered.

This study does not include a replication on a domain independent of the ISIC/PAD-UFES pair, nor a check of generalization beyond the single architecture audited here. Both were considered and are not included: a third-domain replication would require a dataset without documented image overlap with the training data, which was not available within the scope of this study, and a cross-architecture check was not pursued for the same reason. Neither omission is filled by the baseline reference described above, which addresses architecture and dimensionality but not domain independence. The bootstrap check proposed elsewhere to quantify shared estimation noise between condition number and Fisher ratio, both derived from the same fitted precision matrix, was not performed; condition number, not Fisher ratio, is the metric driving the geometric finding in Section~\ref{sec:results-geometry}, which bounds but does not eliminate this as a source of uncertainty in the broader geometry measurement. Checkpoint selection was verified by explicit file path rather than by directory search, but was not independently cross-checked against logged validation curves beyond the single file available per run. Mardia's kurtosis null distribution was calibrated with 200 bootstrap resamples per checkpoint (Section~\ref{sec:geometry-metrics}); the Mahalanobis regularization constant ($\varepsilon=10^{-5}$) and the 16-dimensional representation itself were inherited from the audited classifier rather than varied as part of this study. ``Reliability estimation'' throughout this study refers specifically to out-of-distribution detection AUROC; broader constructs such as calibration or selective prediction are not addressed.

This distinction is particularly relevant when domain-adversarial representations are combined with distance-based or logit-based uncertainty estimation in safety-critical medical AI systems. The same information can remain decodable while becoming largely inaccessible to non-probing reliability estimators.
